% ============================================
% 第1章：机器学会了搭桥，但它不知道自己在搭桥
% ============================================

\chapter{机器学会了搭桥，但它不知道自己在搭桥}

\begin{classroom-scene}
\label{scene:chap1-intro}

{
\noindent \textbf{2025年10月24日，下午5点59分，大连交通大学，桥梁风工程课。}

\vspace{0.5cm}

\noindent\textbf{我：}“好，我们来看一下这个风洞数据。大家用 AI 帮我跑一下这个模态分析，看看前三阶频率是多少。”

\noindent\textbf{学生：}（低头打字，三十秒后）“老师，出来了。一阶0.23，二阶0.58，三阶1.04。”

\noindent\textbf{我：}“这个结果，你们觉得对吗？”

\noindent 教室里安静了五秒。

\noindent\textbf{学生：}“呃\ldots\ldots AI 给的，应该对吧？”

\noindent\textbf{我：}“AI 从来没说过它不会。它只是每次都给出一个答案。”

\noindent 又安静了五秒。

\noindent\textbf{我：}“你们知道吗，AI 搭桥的方式，和你们搭桥的方式，是两种完全不同的东西。它不知道自己在搭桥。它只是算了算，这个像素放在这里，人类看了会说'这是一座桥'。”

\noindent\textbf{学生：}“那\ldots\ldots 它到底是怎么算的？”
}
\end{classroom-scene}

\vspace{0.5cm}

\begin{center}
{\small\color{muted} ——这段对话，发生在我真正开始思考“AI到底怎么学习”的那天。}
\end{center}

\section{AI 的学习：一场盛大的贝叶斯猜谜}

让我们先做一件事。打开 ChatGPT，输入：

\begin{quote}
\texttt{“帮我写一个桥梁工程课程设计任务书，20米简支梁桥，公路一级荷载。”}
\end{quote}

它会给你一份完整的任务书。格式规范，内容合理，甚至附带了参考文献——虽然那些参考文献可能有一半是它编的。

它什么时候学的桥梁工程？它什么时候背过《公路桥涵设计通用规范》？答案是：\textbf{它从来没有“学”过桥梁工程。它只是在海量文本中，发现了“桥梁工程任务书”这个语言模式的各种统计规律。}

这不是学习。这是\textbf{模式匹配}。

\subsection{梯度下降：一条永远在往下走的曲线}

要理解 AI 是怎么“学会”一切的，只需要理解一个词：\textbf{梯度下降}。

想象你站在一座山上，蒙着眼睛，脚下是凹凸不平的地面。你的目标不是看到山顶的风景，而是找到\textbf{最低的那个点}。你能做的只有一件事：每次走一小步，感受脚下的坡度，然后往最陡的下坡方向迈出下一步。重复几百万次。最终，你会停在某个谷底。

这就是 AI 的训练过程。那座山叫“损失函数”——它衡量的是“AI 的预测离正确答案有多远”。梯度下降就是让 AI 不断“往山下走”，每一步都比前一步错得少一点点。

\vspace{0.3cm}

{\small\color{muted} \textbf{技术说明：} 在数学上，梯度下降的更新公式是 $\theta_{t+1} = \theta_t - \eta \nabla L(\theta_t)$，其中 $\theta$ 是模型参数，$\eta$ 是学习率，$\nabla L$ 是损失函数的梯度。把这个公式迭代几十亿次，AI 就“学会”了所有东西。}

\vspace{0.3cm}

\begin{insight}
\label{insight:chap1-1}

\textbf{AI 的学习本质上是一个优化问题。它的目标函数只有一个：让预测误差最小化。}

它不在乎“桥”是什么。它不在乎为什么要建桥。它不在乎那座桥会不会塌。它只是算出了——在它见过的所有文本里——这组参数的组合，能让“预测下一个词”的错误率最低。
\end{insight}

\subsection{上下文窗口：AI 的“记忆”有多长？}

AI 有另一个奇怪的特点：它的“记忆”有一个硬性上限，叫\textbf{上下文窗口}。

你可以把上下文窗口想象成一个只能装固定数量文字的盒子。每次你给 AI 输入新的文字，旧的文字就会被挤出盒子。AI 只能“看见”盒子里的东西——盒子外面的一切，对它来说等于不存在。

2025年，这个盒子大约能装15万字。2026年，能装100万字。但这个数字不管多大，它仍然是一个\textbf{盒子}。AI 没有“经验”。它没有“我记得大一那年，我的结构力学老师说过\ldots\ldots”它没有“上次我做这个的时候，犯了一个错误，这次我不会再犯”。

\begin{insight}
\label{insight:chap1-2}

\textbf{AI 的记忆是“完美但无根”的。}它不会遗忘——只要在上下文窗口里，它就能精确复述。但它也不会“记得”——因为“记得”意味着这个信息曾经改变过你，而你带着这个改变继续生活。
\end{insight}

\subsection{AI 搭桥的方式}

现在我们可以回答那个问题了：AI 是怎么搭桥的？

\textbf{第一步：}它在训练数据中见过几百万个“桥”的文字描述、图片、结构图、规范条文。

\textbf{第二步：}它学会了“桥”这个词在什么语境下出现、和哪些词共现概率高、对应的图片通常有什么样的像素分布。

\textbf{第三步：}当你让它“搭一座桥”，它做的事情是——\textbf{生成一个最符合“桥”这个统计模式的输出}。

它没有计算弯矩。它没有考虑风荷载。它没有检查这根梁会不会被压弯。它只是算出了——在它见过的所有桥里——\textbf{这种形状、这些词、这些像素组合，最像“桥”}。

\begin{center}
\fbox{
\begin{minipage}{0.85\textwidth}
\centering
\textbf{AI 搭桥的逻辑：}

$P(\text{下一个词} \mid \text{上文}) = \text{softmax}(f_\theta(\text{上文}))$

$P(\text{下一个像素} \mid \text{已有像素}) = \text{扩散模型的去噪过程}$

这本质上是\textbf{统计推断}，不是工程设计。
\end{minipage}
}
\end{center}

\section{人的学习：把整座桥的重量扛在肩上}

现在，让我们换一个场景。

\subsection{2025年11月4日，一个学生在搭桥}

\begin{classroom-scene}
\label{scene:chap1-student}

{
\noindent \textbf{2025年11月4日，下午3点28分，桥梁抗风课。}

\vspace{0.5cm}

\noindent 我在电子白板上画了一个矩形。旁边标注：$qL^2/8$。

\noindent\textbf{我：}“简支梁跨中弯矩，$qL^2/8$。这个公式你们大二学过，大三学过，现在大四了，又学一遍。谁能告诉我，这个公式跟搭桥有什么关系？”

\noindent 教室里安静了。

\noindent 然后后排一个学生举手。不是那种“我知道答案”的举手，是那种“我有一个问题但说出来可能很蠢”的举手。

\noindent\textbf{学生：}“老师，我不是不知道这个公式。我是不知道\ldots\ldots 它跟搭桥有什么关系。我知道 $q$ 是荷载，$L$ 是跨度，$qL^2/8$ 是弯矩。但我不知道\ldots\ldots 当我自己去搭桥的时候，这个公式能帮我做什么。”

\noindent 我愣了一下。不是因为这个问题难，而是因为\textbf{这个问题的答案我一直以为大家都懂}。

\noindent\textbf{我：}“好。那我们来搭一座桥。”

\noindent 我打开 SpanCraft，把电脑递给那个学生。他犹豫了一下，点了一下桥面，拖动了一个节点。桥搭好了。他小心翼翼地放了第一辆小车。

\noindent 桥塌了。

\noindent 他眉头皱了一下，加了两个斜杆，改了一下梁高，又放了一辆小车。

\noindent 桥又塌了。

\noindent 这时候发生了一件我没预料到的事。\textbf{他没有抬头看我，没有问“老师怎么办”，而是自己打开了简支梁的那一关，开始调梁高。}

\noindent 他调了三次，桥站住了。他转过头，眼睛里有光：“老师，我好像知道那个 $qL^2/8$ 是什么意思了。\textbf{它说的是——跨中那个地方，弯矩最大，如果你不在那里加高梁，它就会从那里折。}”

\noindent 我什么都没教。他自己发现的。
}
\end{classroom-scene}

\subsection{人学习时，身体在动，手心在出汗}

这个学生经历了一个 AI 永远无法经历的过程。

\textbf{第一，他产生了困惑。}不是“我不知道答案”的困惑，而是“我动手做了，但结果是错的，为什么”的困惑。这种困惑有一个特殊的名字：\textbf{生产性挫折}（productive frustration）。它不是“我不懂”，而是“我做了，但没做成，我需要理解为什么”。

\textbf{第二，他经历了失败。}桥塌了。不是一次，是两次。每次塌，他都能看见塌在哪里。这种反馈是\textbf{即时的、视觉化的、不可否认的}。没有人告诉他“你错了”，是桥自己告诉他的。

\textbf{第三，他自己找到了答案。}他没有被教。他在失败中看见了模式，然后自己去验证了那个模式。那个 $qL^2/8$，从黑板上的一行公式，变成了他指尖上的一根梁。

\textbf{第四，他感到了骄傲。}“老师，我好像知道了。”这句话的语气，不是“我记住了”，而是“我发现了”。这两者之间的区别，就是\textbf{学习}和\textbf{被灌输}之间的区别。

\begin{insight}
\label{insight:chap1-3}

\textbf{人的学习不是模式匹配。人的学习是——经历从不知道到知道的过程，并被这个过程改变。}

AI 的输出是“答案”。人的输出是“我懂了”。AI 不会在输出正确结果时心跳加速。AI 不会在深夜想起自己第一次搭桥成功的那一瞬间，嘴角微微上扬。
\end{insight}

\section{两种学习，两种“理解”}

现在，让我们把 AI 的学习和人的学习放在一起比较。

\begin{table}[H]
\centering
\begin{tabular}{p{0.35\textwidth} p{0.35\textwidth}}
\toprule
\textbf{AI 的学习} & \textbf{人的学习} \\
\midrule
目标：最小化预测误差 & 目标：理解世界，并能在其中行动 \\
\addlinespace
输入：海量数据（文本、图片、代码） & 输入：有限但多模态的体验（视觉、触觉、情感、社交） \\
\addlinespace
记忆：上下文窗口内的完美复制 & 记忆：选择性的、在遗忘中保留本质 \\
\addlinespace
反馈：损失函数数值下降 & 反馈：桥塌了、老师点头了、心里“啊哈”了一声 \\
\addlinespace
错误：统计偏差，需要更多数据 & 错误：\textbf{最有价值的老师}——一次塌桥比一百页规范更难忘 \\
\addlinespace
动机：梯度下降算法驱动 & 动机：\textbf{困惑、好奇、骄傲、不服输、非做不可的理由} \\
\addlinespace
“理解”：能生成正确的输出 & “理解”：能解释、能迁移、能判断、\textbf{能在自己错的时候知道自己错了} \\
\bottomrule
\end{tabular}
\caption{AI的学习 vs 人的学习}
\label{tab:ai-vs-human}
\end{table}

\subsection{AI 不知道“为什么”}

AI 可以告诉你“简支梁跨中弯矩是 $qL^2/8$”。它甚至可以帮你推导这个公式。但它不知道\textbf{为什么这个公式重要}。

它不知道这个公式曾经让多少工程师在深夜反复验算，因为一座桥的成败就在那一个数值上。它不知道塔科马大桥在1940年倒塌的时候，工程师们震惊地发现，他们引以为傲的“静力分析”在风的作用下完全失效了。它不知道那些失败改变了一整个学科。

AI 知道\textbf{事实}。它不知道\textbf{事实的重量}。

\begin{shiyu-note}
\label{note:chap1-1}

\textbf{AI已经会搭桥了，问题是人还想不想过桥。}

“AI 的学习”这个词，本身就是一个陷阱。我们把“输出正确答案”和“理解”混为一谈了。就像我们把“会背乘法表”和“懂数学”混为一谈一样。

AI 的学习本质上是优化。它有一个明确的目标函数——让预测误差最小。它通过梯度下降，迭代几十亿次，找到了一组参数，这组参数能让它在测试集上表现优良。仅此而已。

人的学习不是优化。人的学习是\textbf{转化}。你被一件事情改变了。你不再是你。你看见了以前看不见的东西。你开始用新的方式思考。你会在半夜醒来，突然想通一个问题。你会在失败的时候感到自尊受挫，然后在成功的时候感到一种发自心底的骄傲。这些都不是优化。这些都是\textbf{人的质地}。

AI 可以生成一座完美的桥。但它没有“搭桥”的经验。它没有碰过混凝土，没有感受过钢筋被弯折时的阻力，没有在深夜因为一个参数调不对而摔过键盘。\textbf{它没有“手”。}

所以，AI 时代的教育，真正的问题不是“AI 能替代什么”，而是\textbf{“当 AI 能替代一切，人还剩下什么？”}

人剩下的，恰好是 AI 永远无法拥有的东西：\textbf{困惑、好奇心、挫败感、骄傲、责任感、审美判断、以及“非做不可”的理由。}

这些，才是学习的真正内容。而 AI 的出现，只是帮我们把这些东西从“知识传递”的杂音中清晰地分离了出来。
\end{shiyu-note}

\section{桥的隐喻：两种学习方式之间的鸿沟}

我在教桥梁工程。但说到底，我教的不是桥。

我教的是一个\textbf{隐喻}。桥，是一个连接此岸和彼岸的东西。而此岸和彼岸之间，有三种东西可以跨越：

\begin{enumerate}
  \item \textbf{物理的鸿沟：}河流、峡谷、海峡——这是桥的本职工作。
  \item \textbf{认知的鸿沟：}从“不知道”到“知道”——这是教学的本职工作。
  \item \textbf{存在的鸿沟：}从“我是谁”到“我能成为谁”——这是教育的本职工作。
\end{enumerate}

AI 可以搭第一种桥。它甚至可以在第二种桥的某些部分上帮忙——它可以生成解释、画出图表、给出公式。但第三种桥，它搭不了。

\textbf{因为第三种桥需要一个人相信：你值得被改变。}

而 AI 没有“相信”这个功能。它只有概率分布。

\vspace{0.5cm}

\begin{center}
\fbox{
\begin{minipage}{0.85\textwidth}
\centering
\textbf{这座桥的三层含义：}

\vspace{0.3cm}
\begin{tikzpicture}[scale=0.8]
  % 两岸
  \fill[gray!20] (0,0) rectangle (1,3);
  \fill[gray!20] (9,0) rectangle (10,3);
  % 桥面
  \draw[very thick] (1,2) -- (9,2);
  \draw[very thick] (1,2) .. controls (3,0) and (7,0) .. (9,2);
  % 标注
  \node at (0.5,1.5) [rotate=90] {\small 此岸};
  \node at (9.5,1.5) [rotate=270] {\small 彼岸};
  \node at (5,2.5) {\small 物理的桥：跨越空间};
  \node at (5,0.5) {\small 认知的桥：跨越无知};
  \node at (5,-0.5) {\small\color{blue} \textbf{存在的桥：跨越“我是谁”}};
\end{tikzpicture}
\captionof{figure}{桥的三层含义}
\end{minipage}
}
\end{center}

\vspace{0.5cm}

\section{我们现在在哪里？}

这本书写于一个非常具体的时间点：2025—2026年。

在这个时间点上，AI 已经可以做几乎所有的“知识工作”。它可以写论文、画图、写代码、做视频、生成三维模型、分析数据、写公文\ldots\ldots 它甚至可以通过执业资格考试。但与此同时：

\begin{itemize}
  \item 土木工程被调侃为“天坑专业”，新生报到率在下降。
  \item 课堂上，学生带着 ChatGPT 的答案来上课，老师不知道还能教什么。
  \item 教学创新比赛要求“创新”，但大多数创新只是“加了 AI 工具”。
  \item 铁路局培训教室里，老员工第一次面对 AI 时，那种混合着好奇和恐惧的眼神，我至今记得。
\end{itemize}

我们站在一个很奇怪的十字路口。\textbf{AI 能做一切，但人反而更迷茫了。}不是“不知道该做什么”的迷茫，而是“不知道什么还值得做”的迷茫。

这本书试图回答的，正是这个问题。

\vspace{0.5cm}

\begin{center}
{\large\itshape 在一个机器可以比人更快地学会一切的世界里，}
{\large\itshape 人为什么还要学习？}
{\large\itshape 人为什么还要游戏？}
{\large\itshape 人凭什么还能幸福？}
\end{center}

\vspace{0.5cm}

我的回答，藏在接下来的每一章里。但可以先亮出底牌：

\textbf{人学习，不是因为需要知道什么，而是因为需要成为什么。}

\textbf{人游戏，不是因为想要赢，而是因为想要经历从不会到会的过程，并被这个过程改变。}

\textbf{人幸福，不是因为拥有了什么，而是因为亲手做过什么，并可以指着它说：这是我的。}

桥不会自己存在。学习不会自动发生。幸福不会从天而降。

\textbf{请把手弄脏。}

\vspace{1cm}

\begin{decision-point}
\label{decision:chap1}

\noindent\textbf{这一章，我讲了 AI 的学习和人的学习是两种不同的东西。现在换作你：}

\vspace{0.3cm}

\noindent 打开你的下一节课的 PPT。找到第一页。删掉“课程目标”和“本章重点”。

\noindent 换成一句话。一句话，让你的学生\textbf{想}知道接下来会发生什么。

\noindent 不是“通过本章学习，你将掌握\ldots\ldots”
\noindent 而是“你有没有想过，为什么\ldots\ldots？”

\vspace{0.3cm}

\noindent 就这一句。改完发给我。
\end{decision-point}

\vspace{0.5cm}

\begin{flushright}
{\small\color{muted} 下一章：为什么孩子会蹲在沙坑里搭三个小时的桥——关于游戏的本质。}
\end{flushright}